\documentclass[aspectratio=169,10pt]{beamer}

\usetheme{Madrid}
\usecolortheme{default}
\setbeamertemplate{navigation symbols}{}
\setbeamertemplate{caption}[numbered]
\setbeamertemplate{footline}[frame number]

\usepackage[utf8]{inputenc}
\usepackage[T1]{fontenc}
\usepackage{booktabs}
\usepackage{graphicx}
\usepackage{array}
\usepackage{tabularx}
\usepackage{multirow}
\usepackage{adjustbox}

\definecolor{deepblue}{RGB}{30,75,120}
\definecolor{softgreen}{RGB}{64,130,98}
\definecolor{warmred}{RGB}{174,74,63}
\definecolor{mutedgray}{RGB}{96,106,116}

\setbeamercolor{structure}{fg=deepblue}
\setbeamercolor{frametitle}{fg=white,bg=deepblue}
\setbeamercolor{title}{fg=white,bg=deepblue}
\setbeamercolor{block title}{fg=white,bg=deepblue}
\setbeamercolor{block body}{bg=black!3}

\newcommand{\betterdown}{\(\downarrow\)}
\newcommand{\betterup}{\(\uparrow\)}
\newcommand{\smallnote}[1]{\vspace{0.3em}{\scriptsize\color{mutedgray}#1}}

\title[Results and Discussion]{Results and Discussion}
\subtitle{ACO-based climate regionalization of Serbia}
\author{Aleksandra Ruzic}
\date{}

\begin{document}

\begin{frame}
  \titlepage
\end{frame}

\begin{frame}{Results Narrative}
  \begin{block}{Main question}
    Which ACO configurations produce useful climate regionalizations, and how should they be selected when internal fitness and external validation do not always agree?
  \end{block}

  \vspace{0.5em}
  \begin{columns}[T]
    \begin{column}{0.48\textwidth}
      \textbf{What is compared}
      \begin{itemize}
        \item 3 feature representations
        \item 3 active fitness functions
        \item 7 construction strategies
        \item ACO parameter grid
      \end{itemize}
    \end{column}
    \begin{column}{0.48\textwidth}
      \textbf{How it is judged}
      \begin{itemize}
        \item internal fitness values
        \item vegetation consistency metric
        \item visual map structure
        \item feature and elevation coherence
      \end{itemize}
    \end{column}
  \end{columns}
\end{frame}

\begin{frame}{Reference Regionalizations}
  \begin{columns}[T]
    \begin{column}{0.5\textwidth}
      \textbf{Baselines}
      \begin{itemize}
        \item KMeans: clusters climate variables only.
        \item K\"oppen-Geiger: rule-based climate classes.
        \item Both are evaluated with the same metrics as ACO.
      \end{itemize}
    \end{column}
    \begin{column}{0.5\textwidth}
      \begin{adjustbox}{max width=\textwidth}
      \begin{tabular}{lcccc}
        \toprule
        Method & Centroid & Silh. & SNN & Veg. \\
        \midrule
        KMeans & 0.238 & 0.322 & 0.165 & 0.474 \\
        K\"oppen & 0.368 & 0.415 & 0.106 & 0.443 \\
        \bottomrule
      \end{tabular}
      \end{adjustbox}
      \smallnote{Fitness values are minimized; vegetation metric is maximized.}
    \end{column}
  \end{columns}

  \vspace{0.5em}
  \begin{block}{Baseline interpretation}
    KMeans is stronger on centroid and silhouette fitness; K\"oppen preserves more local graph-neighborhood similarity under SNN.
  \end{block}
\end{frame}

\section{Feature Representation Comparison}

\begin{frame}{Feature Representations: Main Experiment Summary}
  \begin{adjustbox}{max width=\textwidth}
  \begin{tabular}{llccc}
    \toprule
    Fitness objective & Feature set & Best fitness \betterdown & Mean fitness \betterdown & Best vegetation \betterup \\
    \midrule
    \multirow{3}{*}{Centroid}
      & K\"oppen-inspired & 0.232 & 0.461 & \textbf{0.485} \\
      & Raw climate & 0.233 & 0.484 & 0.483 \\
      & PCA & 0.304 & 0.631 & 0.467 \\
    \midrule
    \multirow{3}{*}{Silhouette}
      & K\"oppen-inspired & 0.267 & 0.414 & \textbf{0.483} \\
      & Raw climate & \textbf{0.261} & 0.417 & 0.480 \\
      & PCA & 0.263 & 0.445 & 0.468 \\
    \midrule
    \multirow{3}{*}{SNN}
      & K\"oppen-inspired & 0.083 & 0.349 & \textbf{0.484} \\
      & Raw climate & 0.096 & 0.368 & 0.482 \\
      & PCA & \textbf{0.072} & 0.427 & 0.471 \\
    \bottomrule
  \end{tabular}
  \end{adjustbox}
\end{frame}

\begin{frame}{Feature Representation Takeaways}
  \begin{itemize}
    \item K\"oppen-inspired and raw climate features give similar results across objectives.
    \item PCA behaves differently: it can achieve good internal silhouette/SNN values, but its vegetation metric remains lower.
    \item For centroid fitness, PCA is clearly weaker than the other two representations.
    \item The visual comparison therefore focuses on K\"oppen-inspired features: they are interpretable and competitive.
  \end{itemize}

  \vspace{1em}
  \begin{block}{Interpretation}
    Dimensionality reduction may remove information that is useful for vegetation-based validation, even when it preserves enough structure for some internal objectives.
  \end{block}
\end{frame}

\section{Fitness Influence}

\begin{frame}{Agreement Between Fitness Objectives}
  \begin{columns}[T]
    \begin{column}{0.49\textwidth}
      \textbf{Pearson correlation}
      \begin{adjustbox}{max width=\textwidth}
      \begin{tabular}{lccc}
        \toprule
        & Centroid & Silh. & SNN \\
        \midrule
        Centroid & 1.000 & 0.951 & 0.953 \\
        Silh. & 0.951 & 1.000 & \textbf{0.980} \\
        SNN & 0.953 & \textbf{0.980} & 1.000 \\
        \bottomrule
      \end{tabular}
      \end{adjustbox}
    \end{column}
    \begin{column}{0.49\textwidth}
      \textbf{Spearman correlation}
      \begin{adjustbox}{max width=\textwidth}
      \begin{tabular}{lccc}
        \toprule
        & Centroid & Silh. & SNN \\
        \midrule
        Centroid & 1.000 & 0.906 & 0.901 \\
        Silh. & 0.906 & 1.000 & \textbf{0.967} \\
        SNN & 0.901 & \textbf{0.967} & 1.000 \\
        \bottomrule
      \end{tabular}
      \end{adjustbox}
    \end{column}
  \end{columns}

  \vspace{0.8em}
  \begin{block}{Result}
    Configurations that perform well under one objective often perform well under the others, especially for silhouette and SNN.
  \end{block}
\end{frame}

\begin{frame}{Active Fitness and Vegetation Metric}
  \begin{adjustbox}{max width=0.82\textwidth}
  \begin{tabular}{lcccc}
    \toprule
    Fitness objective & K\"oppen & Raw & PCA & Overall \\
    \midrule
    Centroid & -0.962 & -0.974 & -0.972 & \textbf{-0.970} \\
    Silhouette & -0.828 & -0.916 & -0.928 & -0.886 \\
    SNN & -0.846 & -0.903 & -0.870 & -0.870 \\
    \bottomrule
  \end{tabular}
  \end{adjustbox}

  \smallnote{Pearson correlation. Lower fitness is better; higher vegetation metric is better, so negative correlation indicates agreement.}

  \vspace{1em}
  \begin{itemize}
    \item Fitness and vegetation metric agree over the full configuration set.
    \item The agreement is strongest for centroid fitness.
    \item Full-set correlation does not guarantee that the best-ranked configurations are also best by vegetation.
  \end{itemize}
\end{frame}

\begin{frame}{Top-N Correlation: Agreement Changes With Rank}
  \centering
  \includegraphics[width=0.9\textwidth,height=0.72\textheight,keepaspectratio]{../experiments/results/stage1_topn_analysis/topn_corr_fitness_ecology_combined.png}

  \smallnote{Correlation approaches the full-set value only after many configurations are included; among top-ranked configurations, active fitness is a weaker proxy for vegetation quality.}
\end{frame}

\begin{frame}{Top-Ranked Fitness Is Not Equivalent to Top Vegetation}
  \begin{columns}[T]
    \begin{column}{0.45\textwidth}
      \begin{adjustbox}{max width=\textwidth}
      \begin{tabular}{lc}
        \toprule
        Fitness ranking & Overlap with top-50 vegetation \\
        \midrule
        Centroid & 18 / 50 \\
        Silhouette & 0 / 50 \\
        SNN & 1 / 50 \\
        \bottomrule
      \end{tabular}
      \end{adjustbox}
    \end{column}
    \begin{column}{0.52\textwidth}
      \begin{itemize}
        \item Internal fitness is useful for separating weak from acceptable configurations.
        \item Among strong configurations, it does not reliably select the most vegetation-consistent solution.
        \item This is especially important for silhouette and SNN.
      \end{itemize}
    \end{column}
  \end{columns}

  \vspace{1em}
  \begin{block}{Discussion point}
    External validation is not redundant: it changes which configurations look most attractive.
  \end{block}
\end{frame}

\section{Strategy Influence}

\begin{frame}{Construction Strategy Is the Strongest Algorithmic Factor}
  \begin{adjustbox}{max width=\textwidth}
  \begin{tabular}{lcccc}
    \toprule
    Construction strategy & Centroid \betterdown & Silh. \betterdown & SNN \betterdown & Veg. \betterup \\
    \midrule
    ACO Centroid & \textbf{0.318} & \textbf{0.325} & \textbf{0.203} & \textbf{0.461} \\
    ACO Centroid + Refinement & 0.391 & 0.392 & 0.305 & 0.455 \\
    Multi-colony Centroid Graph Spatial & 0.418 & 0.354 & 0.264 & 0.441 \\
    ACO Graph Local & 0.594 & 0.442 & 0.375 & 0.427 \\
    Multi-colony Cohesion Graph Spatial & 0.619 & 0.462 & 0.423 & 0.425 \\
    ACO Cohesion + Refinement & 0.662 & 0.500 & 0.547 & 0.422 \\
    ACO Cohesion & 0.677 & 0.503 & 0.550 & 0.419 \\
    \bottomrule
  \end{tabular}
  \end{adjustbox}

  \smallnote{Mean fitness by active objective; vegetation metric averaged across objectives.}
\end{frame}

\begin{frame}{Strategy Winners Depend on the Objective}
  \begin{adjustbox}{max width=\textwidth}
  \begin{tabular}{lll}
    \toprule
    Ranking criterion & Dominant top strategies & Interpretation \\
    \midrule
    Centroid fitness & ACO Centroid & strong compactness around climate centroids \\
    Silhouette fitness & mostly ACO Centroid & centroid construction remains competitive \\
    SNN fitness & multi-colony / graph strategies appear & graph-local structure matters more \\
    Vegetation metric & ACO Centroid, Centroid + Refinement & external validation favors centroid family \\
    \bottomrule
  \end{tabular}
  \end{adjustbox}

  \vspace{1em}
  \begin{block}{Answer to the main concern}
    Centroid construction does not win only because the centroid fitness favors it. It remains strong under silhouette fitness and in vegetation validation, while SNN makes graph-aware strategies more competitive.
  \end{block}
\end{frame}

\begin{frame}{Top Configurations: What Repeats?}
  \begin{columns}[T]
    \begin{column}{0.5\textwidth}
      \textbf{Repeated pattern}
      \begin{itemize}
        \item Top centroid and silhouette configurations are dominated by ACO Centroid.
        \item Top SNN configurations are more mixed.
        \item Top vegetation configurations remain centroid or centroid-refinement based.
      \end{itemize}
    \end{column}
    \begin{column}{0.5\textwidth}
      \textbf{Interpretation}
      \begin{itemize}
        \item Strategy quality should be interpreted together with the active objective.
        \item Centroid-based methods are the most robust family in this experimental grid.
        \item SNN is useful because it exposes alternatives that compactness-based objectives may hide.
      \end{itemize}
    \end{column}
  \end{columns}
\end{frame}

\section{Parameter Influence}

\begin{frame}{Parameter Influence Summary}
  \begin{adjustbox}{max width=\textwidth}
  \begin{tabular}{llcccc}
    \toprule
    Parameter & Value & Centroid \betterdown & Silh. \betterdown & SNN \betterdown & Veg. \betterup \\
    \midrule
    \multirow{3}{*}{Pheromone update}
      & All ants & \textbf{0.502} & \textbf{0.417} & \textbf{0.352} & \textbf{0.439} \\
      & Best ant & 0.545 & 0.431 & 0.402 & 0.434 \\
      & Quality weighted & 0.530 & 0.428 & 0.390 & 0.435 \\
    \midrule
    \multirow{3}{*}{\(\alpha\)}
      & 0.75 & 0.549 & 0.432 & 0.400 & 0.434 \\
      & 1.00 & \textbf{0.506} & \textbf{0.412} & \textbf{0.356} & \textbf{0.438} \\
      & 2.00 & 0.522 & 0.432 & 0.387 & 0.435 \\
    \midrule
    \multirow{2}{*}{\(\beta\)}
      & 1 & 0.684 & 0.508 & 0.547 & 0.423 \\
      & 3 & \textbf{0.367} & \textbf{0.343} & \textbf{0.215} & \textbf{0.449} \\
    \midrule
    \multirow{2}{*}{Evaporation}
      & 0.02 & 0.527 & 0.426 & 0.382 & 0.436 \\
      & 0.10 & \textbf{0.524} & \textbf{0.425} & \textbf{0.380} & \textbf{0.436} \\
    \bottomrule
  \end{tabular}
  \end{adjustbox}
\end{frame}

\begin{frame}{Main Parameter Result: Heuristic Weight}
  \begin{columns}[T]
    \begin{column}{0.48\textwidth}
      \textbf{\(\beta = 3\) improves all objectives}
      \begin{itemize}
        \item Centroid mean: 0.684 \(\rightarrow\) 0.367
        \item Silhouette mean: 0.508 \(\rightarrow\) 0.343
        \item SNN mean: 0.547 \(\rightarrow\) 0.215
        \item Vegetation mean: 0.423 \(\rightarrow\) 0.449
      \end{itemize}
    \end{column}
    \begin{column}{0.48\textwidth}
      \textbf{Other parameters}
      \begin{itemize}
        \item \(\alpha=1\) is best on average, but with smaller gaps.
        \item All-ants update is strongest on average.
        \item Evaporation has almost no visible effect in the tested range.
      \end{itemize}
    \end{column}
  \end{columns}

  \vspace{1em}
  \begin{block}{Interpretation}
    Heuristic information has a stronger influence on solution quality than the tested pheromone or evaporation settings.
  \end{block}
\end{frame}

\section{Visual and Physical Interpretation}

\begin{frame}{Reference Maps}
  \centering
  \includegraphics[width=0.86\textwidth,height=0.76\textheight,keepaspectratio]{../experiments/results/report_comparison_outputs/reference_regionalizations_koppen_features.png}

  \smallnote{KMeans groups by climate-variable similarity; K\"oppen-Geiger uses predefined climate rules.}
\end{frame}

\begin{frame}{Selected ACO Regionalizations}
  \centering
  \includegraphics[width=0.96\textwidth,height=0.78\textheight,keepaspectratio]{../experiments/results/report_comparison_outputs/selected_aco_regionalizations_koppen_features.png}

  \smallnote{One solution is selected for each active fitness function, plus one solution selected by vegetation metric.}
\end{frame}

\begin{frame}{Cluster Size and Dominant Regions}
  \begin{adjustbox}{max width=0.86\textwidth}
  \begin{tabular}{lc}
    \toprule
    Regionalization & Largest cluster share \\
    \midrule
    KMeans & 0.352 \\
    K\"oppen-Geiger & 0.614 \\
    Best centroid fitness & 0.310 \\
    Best silhouette fitness & 0.739 \\
    Best SNN fitness & 0.800 \\
    Best vegetation metric & 0.302 \\
    \bottomrule
  \end{tabular}
  \end{adjustbox}

  \vspace{1em}
  \begin{itemize}
    \item Silhouette and SNN selected maps are dominated by one large region.
    \item Centroid and vegetation-selected ACO maps are more balanced.
    \item This motivates checking whether visual dominance corresponds to weaker internal physical coherence.
  \end{itemize}
\end{frame}

\begin{frame}{Within-Cluster Climate Feature Variation}
  \begin{adjustbox}{max width=\textwidth}
  \begin{tabular}{lcc}
    \toprule
    Regionalization & Mean temperature ratio \betterdown & Total precipitation ratio \betterdown \\
    \midrule
    KMeans & \textbf{0.431} & \textbf{0.393} \\
    K\"oppen-Geiger & 0.532 & 0.724 \\
    Best centroid fitness & 0.459 & 0.414 \\
    Best silhouette fitness & 0.715 & 0.656 \\
    Best SNN fitness & 0.682 & 0.663 \\
    Best vegetation metric & 0.461 & 0.431 \\
    \bottomrule
  \end{tabular}
  \end{adjustbox}

  \smallnote{Ratios are pooled within-cluster standard deviations divided by the global standard deviation of that feature.}

  \vspace{0.7em}
  \begin{block}{Result}
    KMeans, best centroid-fitness and best vegetation-metric solutions have lower internal climate-feature variation. The selected silhouette and SNN solutions are less coherent for temperature.
  \end{block}
\end{frame}

\begin{frame}{Elevation as Independent Validation}
  \begin{columns}[T]
    \begin{column}{0.5\textwidth}
      \begin{adjustbox}{max width=\textwidth}
      \begin{tabular}{lc}
        \toprule
        Regionalization & Pooled elevation std. (m) \betterdown \\
        \midrule
        KMeans & \textbf{202.1} \\
        K\"oppen-Geiger & 239.4 \\
        Best centroid fitness & 209.0 \\
        Best silhouette fitness & 302.2 \\
        Best SNN fitness & 293.3 \\
        Best vegetation metric & 210.1 \\
        \bottomrule
      \end{tabular}
      \end{adjustbox}
    \end{column}
    \begin{column}{0.47\textwidth}
      \begin{itemize}
        \item Elevation is not used during clustering.
        \item Centroid and vegetation-selected ACO maps are close to KMeans.
        \item Silhouette and SNN selected maps have much higher elevation variation.
      \end{itemize}
    \end{column}
  \end{columns}
\end{frame}

\begin{frame}{High-Elevation Concentration}
  \begin{adjustbox}{max width=0.82\textwidth}
  \begin{tabular}{lcc}
    \toprule
    Method & Share of top cluster above 1000 m & Captured pixels above 1000 m \\
    \midrule
    K\"oppen-Geiger & 31.8\% & 95.5\% \\
    Best vegetation metric & \textbf{57.4\%} & 84.7\% \\
    \bottomrule
  \end{tabular}
  \end{adjustbox}

  \vspace{1em}
  \begin{itemize}
    \item K\"oppen captures most high-elevation pixels, but places them in a broader mountain-containing cluster.
    \item The best vegetation ACO solution forms a more concentrated high-elevation cluster.
    \item This supports the idea that external validation can reveal physically meaningful differences between maps.
  \end{itemize}
\end{frame}

\section{Multi-Seed Validation}

\begin{frame}{Multi-Seed Validation}
  \begin{itemize}
    \item Selected configurations were rerun with five seeds.
    \item The goal was to check whether the strongest main-experiment patterns survive additional stochastic variation.
    \item Validation supports strategy-family conclusions more than exact parameter rankings.
  \end{itemize}

  \vspace{0.7em}
  \begin{adjustbox}{max width=\textwidth}
  \begin{tabular}{lll}
    \toprule
    Active objective & Validated pattern & Main interpretation \\
    \midrule
    Centroid & ACO Centroid remains best across feature sets & stable centroid-family result \\
    Silhouette & mostly ACO Centroid; one multi-colony winner & centroid remains strong \\
    SNN & mixed winners across feature sets & graph/cohesion strategies become competitive \\
    Vegetation & centroid or centroid-refinement winners & external metric favors centroid family \\
    \bottomrule
  \end{tabular}
  \end{adjustbox}
\end{frame}

\begin{frame}{Statistical Caution}
  \begin{itemize}
    \item Pairwise paired t-tests often indicated differences between selected configurations.
    \item Wilcoxon signed-rank tests did not produce pairwise differences with \(p<0.05\).
    \item With five paired seeds, Wilcoxon has very low resolution; in many comparisons the smallest attainable two-sided \(p\)-value is 0.0625.
  \end{itemize}

  \vspace{1em}
  \begin{block}{Validated conclusion}
    The results justify broad qualitative claims about strategy families, but not strong claims that one exact top configuration is statistically superior to another.
  \end{block}
\end{frame}

\section{Conclusions}

\begin{frame}{Main Conclusions}
  \begin{itemize}
    \item ACO is a useful framework for producing alternative climate regionalizations of Serbia.
    \item Centroid-based construction is the most robust strategy family in the tested grid.
    \item The selected fitness objective matters: SNN makes graph-aware and multi-colony strategies more competitive.
    \item K\"oppen-inspired and raw features are similarly strong; PCA is less convincing under vegetation validation.
    \item The vegetation metric is not redundant with internal fitness and can change which solution is preferred.
    \item Elevation and within-cluster feature variation support the selected centroid/vegetation solutions as physically coherent.
  \end{itemize}
\end{frame}

\begin{frame}{Final Discussion Point}
  \begin{block}{Practical recommendation}
    Use internal fitness to remove weak configurations, then select among strong candidates using external validation and map interpretation.
  \end{block}

  \vspace{1em}
  \begin{itemize}
    \item If compact climate regions are the priority: ACO Centroid is the safest choice.
    \item If graph-local structure is the priority: SNN objectives deserve separate interpretation.
    \item If ecological plausibility is the priority: vegetation and elevation checks should be part of final selection.
  \end{itemize}
\end{frame}

\end{document}
